\chapter{Producing PDS and FFT Products}

The first GHATS step for most analyses is to turn mission event data, binned
data, or an external light curve into a native GHATS timing product.  There
are two closely related products:

\begin{itemize}
\item a \texttt{.pds} file, containing one power spectrum for each accepted
  uninterrupted data stretch;
\item a \texttt{.fft} file, containing the corresponding complex Fourier
  amplitudes.
\end{itemize}

The two files are made by the same family of mission-specific production
routines.  They use the same transform lengths, time binning, channel
selection, GTI handling, and GHATS header conventions.  The difference is only
what is stored after each transform has been computed: powers for a
\texttt{.pds} file, or complex Fourier amplitudes for a \texttt{.fft} file.

\begin{ghatsnote}
Use a \texttt{.pds} file when the next step is ordinary power-spectral
analysis.  Use a \texttt{.fft} file when the next step needs phase information,
for example cross spectra, lags, coherence, or bispectral products.
\end{ghatsnote}

\section{Common Calling Sequence}

Most mission routines have the same basic interface:

\begin{ghatsconsole}
Ghats> gh_xte
Ghats> gh_xte, '#parameters.par'
Ghats> gh_xte, infile, outtype, channels, treb, npds, outfile
\end{ghatsconsole}

The first form runs interactively.  The second form reads a parameter file.
The third form gives the main parameters directly on the IDL/GDL command line.
For a different mission, replace \routine{gh_xte} by the appropriate routine,
for instance \routine{gh_nicer}, \routine{gh_ep}, or \routine{gh_nustar}.

\begin{description}
\item[\param{infile}] Input science file, \texttt{@} metafile, or
  \texttt{@@} meta-metafile.  A metafile contains one input filename per line.
\item[\param{outtype}] Output type.  Use \texttt{'POWER'} for a
  \texttt{.pds} file or \texttt{'FFT'} for a \texttt{.fft} file.
\item[\param{channels}] Two-element channel interval, for example
  \texttt{[0,35]}.  The meaning of the channel is mission dependent: RXTE/PCA
  uses PCA channels, while event-list missions normally use PI or PHA channels.
\item[\param{treb}] Integer rebinning factor in time before computing the
  transforms.  A value of 1 keeps the native time resolution.
\item[\param{npds}] Transform size.  If this is an integer or long integer,
  it is interpreted as the number of rebinned light-curve bins per transform.
  If this is a real number, it is interpreted as a duration in seconds.
\item[\param{outfile}] Output filename, normally ending in \texttt{.pds} or
  \texttt{.fft}.
\end{description}

For example:

\begin{ghatsconsole}
Ghats> gh_xte, '@event.lis', 'POWER', [0,249], 1, 8192, 'fullband.pds'
Ghats> gh_xte, '@event.lis', 'FFT',   [0,13],  1, 8192, 'soft.fft'
Ghats> gh_xte, '@event.lis', 'FFT',   [14,249],1, 8192, 'hard.fft'
\end{ghatsconsole}

The two \texttt{.fft} files in the last two examples can then be used for
cross-spectral work, provided they contain the same accepted transforms in the
same order.

\section{Interactive and Parameter-File Modes}

When called without arguments, a production routine asks the user for the same
choices that appear in the command-line form: input file, output type, channel
range, time rebinning, transform length, and output filename.  This is often
the safest way to make a first product from an unfamiliar observation, because
the routine prints the channel information, time resolution, and final
transform summary as it goes.

A parameter file gives the same six arguments to the routine.  The command is
called with a leading \texttt{\#}:

\begin{ghatsconsole}
Ghats> gh_xte, '#mysource.par'
\end{ghatsconsole}

This is useful when the same production command must be repeated exactly, or
when several energy bands are being produced from the same input data.  The
parameter-file route should preserve the same choices as the explicit
command-line route; if two products will later be compared transform by
transform, make the choices identical except for the channel range.

\section{Input Lists}

The input filename may point to a single file or to a list.  GHATS uses a
leading \texttt{@} to identify a metafile:

\begin{ghatsconsole}
$ ls *evt* > event.lis
Ghats> gh_nicer, '@event.lis', 'POWER', [30,1200], 1, 4096, 'nicer.pds'
\end{ghatsconsole}

A leading \texttt{@@} identifies a meta-metafile, that is, a file containing
names of metafiles.  This convention is inherited from the RXTE routines and
is useful when one observation contains several groups of homogeneous input
files.

Files grouped in one list should be compatible: same data mode or event-file
layout, same time resolution, same relevant detector setup, and the same
energy-channel convention.  GTIs and gaps may differ from file to file; the
production routine constructs transforms only from accepted continuous
stretches.

\section{Choosing the Transform Length}

The transform length determines the frequency grid.  If \(N_{\rm pds}\) rebinned
bins enter each transform and the rebinned time resolution is \(\Delta t\),
then the transform duration is

\[
  T_{\rm seg} = N_{\rm pds}\,\Delta t .
\]

The frequency resolution is \(1/T_{\rm seg}\), and the Nyquist frequency is
\(1/(2\Delta t)\).  In GHATS files, the DC bin is not stored and the positive
Fourier bins through the Nyquist bin are stored.

The production routines also accept the transform length in seconds.  For
example, if the data after rebinning have a time resolution of \(1/8192\) s,
then these two calls request the same transform duration:

\begin{ghatsconsole}
Ghats> gh_xte, '@event.lis', 'POWER', [0,249], 1, 262144L, 'event.pds'
Ghats> gh_xte, '@event.lis', 'POWER', [0,249], 1, 32.0,    'event.pds'
\end{ghatsconsole}

\begin{ghatsnote}
The number of points per transform should be a power of two.  Long transform
lengths give finer frequency resolution but require longer uninterrupted data
stretches.  Shorter transform lengths keep more stretches when the observation
contains many gaps.
\end{ghatsnote}

\section{Channels and Standard Bands}

The main \param{channels} argument selects the channel range accumulated into
the timing product itself.  Most production routines also accept
\keyword{BANDS}.  This keyword defines three auxiliary channel bands stored in
the GHATS product and used later by routines such as \routine{gh_colors}.

\begin{ghatsconsole}
Ghats> gh_xte, '@event.lis', 'POWER', [0,249], 1, 4096, 'event.pds', $
       bands=[3,6,7,14,15,20]
\end{ghatsconsole}

The six numbers define three intervals: \texttt{3--6}, \texttt{7--14}, and
\texttt{15--20}.  They are instrument channels, not automatically calibrated
physical energy bands.  For missions such as Einstein Probe, the default
split is a technical PI-channel split unless the user supplies a more
appropriate one.

\section{GTIs, Barycentered Times, and Sliding Windows}

Several common keywords control how the input data are segmented:

\begin{description}
\item[\keyword{GTI}] Intersect the data GTIs with user-supplied GTIs.  The
  user GTI file may be an ASCII file with start and stop times or a standard
  GTI FITS file, depending on the mission routine.
\item[\keyword{BARY}] Read barycentered times from the appropriate
  barycentered time column, normally \texttt{BARYTIME}.  The user is
  responsible for ensuring that the column exists and is appropriate.
\item[\keyword{SLIDING}] Use overlapping transforms.  If
  \texttt{SLIDING=4}, consecutive transforms are shifted by one quarter of
  the transform duration.  This can be useful diagnostically, but the averaged
  PDS obtained from overlapping stretches does not have the same simple
  statistical interpretation as an average of independent transforms.
\end{description}

\section{Window Functions}

The production routines accept \keyword{WIND} and \param{WPAR}, which are
passed to the GHATS windowing machinery.  If no window is requested, the
default is a boxcar window.  A non-boxcar window changes the effective
weighting of the light curve before the Fourier transform and should therefore
be recorded when reporting a product.

\section{Mission-Specific Routines}

\begin{longtable}{>{\raggedright\arraybackslash}p{0.24\textwidth}>{\raggedright\arraybackslash}p{0.20\textwidth}>{\raggedright\arraybackslash}p{0.42\textwidth}}
\toprule
Routine & Mission & Purpose \\
\midrule
\endhead
\routine{gh_xte} & RXTE/PCA & Produce GHATS PDS or FFT files from RXTE/PCA data. \\
\routine{gh_nicer} & NICER & Produce GHATS PDS or FFT files from NICER event data. \\
\routine{gh_ep} & Einstein Probe & Produce GHATS PDS or FFT files from Einstein Probe event data. \\
\routine{gh_nustar} & NuSTAR & Produce GHATS PDS or FFT files from NuSTAR event data. \\
\routine{gh_xmm} & XMM & Produce GHATS PDS or FFT files from XMM event data. \\
\routine{gh_swift} & Swift/XRT & Produce GHATS PDS or FFT files from Swift/XRT event data. \\
\routine{gh_ixpe} & IXPE & Produce GHATS PDS or FFT files from IXPE event data. \\
\routine{gh_laxpc} & AstroSat/LAXPC & Produce GHATS PDS or FFT files from AstroSat/LAXPC data. \\
\routine{gh_laxpc_misra} & AstroSat/LAXPC & Produce GHATS PDS or FFT files from AstroSat/LAXPC event data. \\
\routine{gh_czti} & AstroSat/CZTI & Produce GHATS PDS or FFT files from AstroSat/CZTI event data. \\
\routine{gh_huiyan} & Huiyan/HXMT & Produce GHATS PDS or FFT files from Huiyan/HXMT data. \\
\routine{gh_huiyan_le} & Huiyan/HXMT LE & Produce GHATS PDS or FFT files from Huiyan/HXMT LE data. \\
\routine{gh_huiyan_me} & Huiyan/HXMT ME & Produce GHATS PDS or FFT files from Huiyan/HXMT ME data. \\
\routine{gh_huiyan_he} & Huiyan/HXMT HE & Produce GHATS PDS or FFT files from Huiyan/HXMT HE data. \\
\routine{gh_ascii_lcfft} & Contributed & Produce GHATS PDS or FFT files from an ASCII light curve. \\
\bottomrule
\end{longtable}

Each routine can be queried from GHATS:

\begin{ghatsconsole}
Ghats> gh_xte, /help
Ghats> gh_nicer, /help
Ghats> gh_ep, /help
\end{ghatsconsole}

\section{RXTE/PCA}

\routine{gh_xte} is the original GHATS production routine and supports RXTE
binned, single-bit, event, and GoodXenon-style inputs.  It also supports
\keyword{TURBO}, which uses the faster RADPS path where available, and
\keyword{BANDS} for the three auxiliary Standard-2-like rate bands.

\begin{ghatsconsole}
Ghats> gh_xte, '@event.lis', 'POWER', [0,249], 1, 4096, 'rxte_full.pds'
Ghats> gh_xte, '@event.lis', 'FFT',   [0,13],  1, 4096, 'rxte_soft.fft'
\end{ghatsconsole}

For RXTE/PCA PDS products, GHATS can store and later use PCA-specific
information relevant to Poisson-noise estimates.  The standard PDS values
written by the usual production path are in Leahy normalization.  The
\texttt{.fft} products store raw complex Fourier amplitudes; normalizations
are applied by later analysis routines.

\section{Event-List Missions}

Several newer routines follow the same event-list template:
\routine{gh_nicer}, \routine{gh_ep}, \routine{gh_nustar}, \routine{gh_xmm},
\routine{gh_swift}, and \routine{gh_ixpe}.  Their common command-line form is:

\begin{ghatsconsole}
Ghats> gh_nicer, infile, outtype, channels, treb, npds, outfile
\end{ghatsconsole}

The most important mission-dependent details are the event time column, the
channel column, the native time system, and the default auxiliary bands.  The
routine help should be checked before first use:

\begin{ghatsconsole}
Ghats> gh_nicer, /help
Ghats> gh_ep, /help
\end{ghatsconsole}

Examples:

\begin{ghatsconsole}
Ghats> gh_nicer, 'ni_clean.evt', 'FFT', [30,1200], 10000, 32768, 'ni.fft'
Ghats> gh_ep, 'fxt_clean.fits', 'POWER', [0,1023], 1, 4096, 'fxt.pds'
\end{ghatsconsole}

\begin{ghatsnote}
\routine{gh_nicer} has the keyword \keyword{NONOISYDET} for NICER-specific
noisy-detector filtering.  This is not a generic event-list option and is not
used by \routine{gh_ep}.
\end{ghatsnote}

\section{AstroSat and Huiyan/HXMT}

AstroSat and Huiyan/HXMT have mission-specific front ends because their data
layouts and detector selections differ from the simpler event-list template.
The output products, however, are ordinary GHATS \texttt{.pds} and
\texttt{.fft} files and can be read by the same downstream routines.

For AstroSat:

\begin{ghatsconsole}
Ghats> gh_laxpc, infile, outtype, channels, treb, npds, outfile
Ghats> gh_laxpc_misra, infile, outtype, channels, units, layers, treb, npds, outfile
Ghats> gh_czti, infile, outtype, channels, treb, npds, outfile
\end{ghatsconsole}

For Huiyan/HXMT:

\begin{ghatsconsole}
Ghats> gh_huiyan, infile, outtype, channels, treb, npds, outfile
Ghats> gh_huiyan_le, infile, outtype, channels, treb, npds, outfile
Ghats> gh_huiyan_me, infile, outtype, channels, treb, npds, outfile
Ghats> gh_huiyan_he, infile, outtype, channels, treb, npds, outfile
\end{ghatsconsole}

Use the routine-specific \keyword{HELP} output for the exact channel,
detector, unit, and layer arguments.

\section{ASCII Light Curves}

\routine{gh_ascii_lcfft} is a contributed routine for producing a GHATS
\texttt{.pds} or \texttt{.fft} file from an ASCII light curve with columns

\begin{ghatsconsole}
time(s)   counts_or_rate   error
\end{ghatsconsole}

The error column is accepted for compatibility with common light-curve
formats, but is not used in the Fourier calculation.  By default, the second
column is interpreted as counts per original light-curve bin.  If it is a
count rate, use \keyword{RATE}.

\begin{ghatsconsole}
Ghats> gh_ascii_lcfft, 'lc.txt', 'POWER', 1, 4096, 'lc.pds', dt=0.016d0
Ghats> gh_ascii_lcfft, 'lc_rate.txt', 'FFT', 1, 8192, 'lc.fft', $
       dt=0.001d0, /rate
\end{ghatsconsole}

The routine detects gaps in the input time column and does not Fourier
transform across them.  Rows with non-finite times or rates are treated as
boundaries rather than interpolated over.

\section{After Production}

After making a product, inspect it before using it in later analysis:

\begin{ghatsconsole}
Ghats> gh_info, 'event.pds'
Ghats> ghats_all, 'event.pds'
\end{ghatsconsole}

For \texttt{.fft} files, \routine{gh_info} is the appropriate header check.
\routine{ghats_all} can print header information from FFT products, but its
summary plots are designed for PDS files.

The main downstream readers are:

\begin{itemize}
\item \routine{ghx}, to read and average \texttt{.pds} products;
\item \routine{ghx_fft}, to read and average \texttt{.fft} products;
\item \routine{gh_licu}, to extract the stored light curve;
\item \routine{gh_cs} and related cross-spectral routines, to work from
  matching FFT products.
\end{itemize}

For the exact binary layout, frequency convention, FFT sign convention, and
normalization details, see Appendix~\ref{app:file-formats}.
